\chapter{Credible Empirical Work}\label{ch:b3:credible-empirical-work}

The first five chapters built designs that can identify a causal effect:
instruments, panels, differences in differences, discontinuities. A design is
a promise about one study. Economics, like every empirical field, learns from
many studies, and what it learns depends on which studies are run, which
analyses are reported and which results are published. If significant results
are published and others are not, the published record can be wrong even when
every study in it is honest. This chapter measures how large that problem is
in economics, shows the tools that detect it and the practices that limit it,
and ends with the scale on which this volume grades evidence from here on.

\begin{openingchart}{Simulated: the absolute $z$-statistics of 6{,}000
studies, each with its own true effect and sample size (light bars), and of the
2{,}663 that are published (dark bars) when every study with $|z| > 1.96$ is
published and the others one time in four. Bins of width 0.245, so that 1.96 is
a bin edge.}{Simulated by the authors (seed 1996)}
\begin{tikzpicture}
\begin{axis}[omtime, width=0.8\linewidth, height=0.42\linewidth, xmin=0, xmax=5.9,
  ymin=0, ymax=800, xlabel={Absolute $z$-statistic}, ylabel={Studies per bin},
  xtick={0,1,1.96,3,4,5}, xticklabels={0,1,1.96,3,4,5}, legend pos=north east]
  \addplot[ybar, bar shift=0pt, bar width=7.5pt, area legend, fill=omDef!20, draw=omDef!50]
    table[x=z, y=all] {figdata/bachelor-3/06-credible-empirical-work/out/filter_hist.csv};
  \addplot[ybar, bar shift=0pt, bar width=7.5pt, area legend, fill=omDef!75, draw=omDef]
    table[x=z, y=published] {figdata/bachelor-3/06-credible-empirical-work/out/filter_hist.csv};
  \draw[guide] (axis cs:1.96,0) -- (axis cs:1.96,800);
  \legend{all studies, published}
\end{axis}
\end{tikzpicture}
\end{openingchart}

The light bars fall smoothly: most studies are small, most true effects are
modest, and $z$-statistics near zero are the most common. The dark bars do not.
Below 1.96 they are a quarter of the light ones; above it they are all of
them, and the published record has a cliff exactly where a significance test
draws its line. In the simulation, 27~\% of all studies are significant, and
61~\% of the published ones. % figdata: filter_meta
Real collections of published test statistics look like this. In about 50{,}000
tests published in three leading economics journals, Brodeur and coauthors find
the distribution of $p$-values shaped like a two-humped camel, with $p$-values
between 0.25 and 0.10 missing. % ledger: b3ce:blsz.camel
The rest of the chapter asks what such a filter does to estimates, how to see
it, and what to do about it.

\begin{recall}
The Year 1 volume defined tests, $p$-values and the power of a test, and
showed what happens when many tests are run at once. The Year 2 volume named
\emph{specification search}: trying several specifications and reporting the
ones that work. Chapters~\ref{ch:b3:instrumental-variables}
to~\ref{ch:b3:regression-discontinuity} supplied the designs whose results are
graded at the end of this chapter.
\end{recall}

\begin{notation}[Symbols of this chapter]\label{not:b3:credible-empirical-work:symbols}
$\beta$ is a true effect and $\hat\beta$ an estimate of it, with standard error
$s$; $z = \hat\beta/s$; $c = 1.96$ is the critical value. $\phi$ and $\Phi$ are the
standard normal density and distribution function, and $\lambda(a) =
\phi(a)/(1 - \Phi(a))$. In \omterm{def:b3:credible-empirical-work:meta-analysis}{meta-analysis}, study $i$ reports $\hat\beta_i$ with
standard error $s_i$, and $\tau^2$ is the variance of true effects across
studies (not a tax rate).
\end{notation}

\section{How results get published}\label{sec:b3:credible-empirical-work:publication}

A study is designed, run, analysed and written up; then a journal decides. At
each step a significant result has an advantage: researchers are more likely to
write it up, referees more likely to find it interesting, editors more likely
to accept it. None of this requires dishonesty.

\begin{definition}[Publication bias]\label{def:b3:credible-empirical-work:publication-bias}
\index{publication bias}
\emph{Publication bias} is the difference between the average published
estimate of an effect and the effect itself, created when whether a result is
published depends on its sign, its size or its statistical significance.
\end{definition}

\begin{proposition}[The significance filter inflates published estimates]\label{prop:b3:credible-empirical-work:filter}
Let $\hat\beta \sim N(\beta, s^2)$ with $\beta > 0$, and let the estimate be
published only if $\hat\beta/s > c$. Write $a = c - \beta/s$. Then the
probability of publication is the power $1 - \Phi(a)$, and
\[
  \E[\hat\beta \mid \text{published}] = \beta + s\,\lambda(a) > \beta .
\]
\end{proposition}

\begin{proof}
Write $\hat\beta = \beta + s u$ with $u \sim N(0, 1)$. The estimate is published
when $u > a$, which has probability $1 - \Phi(a)$. Then $\E[\hat\beta \mid u > a]
= \beta + s\,\E[u \mid u > a]$, and
\[
  \E[u \mid u > a] = \frac{1}{1 - \Phi(a)} \int_a^\infty u\,\phi(u)\, du
  = \frac{\bigl[-\phi(u)\bigr]_a^\infty}{1 - \Phi(a)} = \frac{\phi(a)}{1 - \Phi(a)},
\]
using $\phi'(u) = -u\phi(u)$. Since $\phi(a) > 0$, the term is positive.
\end{proof}

The size of the bias is $s\,\lambda(a)$: it grows with the standard error, so the
smallest studies are the most exaggerated, and it is largest when power is low.

\begin{model}{The significance filter}
\assumes Studies estimate an effect without bias; a study is published if its
estimate is significant, and otherwise with a lower probability.
\predicts A heap of published test statistics just above the critical value
and a gap just below it; published estimates larger than the true effect, by
more in small studies; estimates that rise with their standard errors across a
literature.
\tested Collections of published test statistics in economics show the gap below
the critical value; across 159 literatures, most reported effects are larger
than those of the adequately powered studies of the same literature.
\breaks It is not the only reason for the heap: specification search produces
it too, within studies, and the two are hard to tell apart in published numbers.
Where null results are wanted --- a \omterm{def:b3:credible-empirical-work:replication}{replication}, a policy that is claimed not to
work --- the filter can run the other way.
\end{model}

\section{Specification search and the garden of forking paths}\label{sec:b3:credible-empirical-work:forking}

Even before a journal decides, the analysis has made choices. Which outcome?
Which controls? Which sample? What to do with outliers? Each choice is often
defensible, and a researcher who makes them after seeing the data can, without
intending to, follow the path that gives the clearest result.

\begin{proposition}[Many independent tests]\label{prop:b3:credible-empirical-work:many}
If $k$ independent tests are run at the 5~\% level and every null hypothesis is
true, the probability that at least one is significant is $1 - 0.95^k$.
\end{proposition}

\begin{proof}
Each test is not significant with probability 0.95; by independence, all $k$ are
not significant with probability $0.95^k$; at least one is significant with the
complementary probability.
\end{proof}

For $k = 16$ the formula gives 0.56. The paths of one analysis are not
independent --- they use the same data --- so the true figure is lower, but it
is far above 5~\%. The diagram shows sixteen paths for one simple question, and
\cref{lst:b3:credible-empirical-work:forking} walks them on data in which the
treatment does nothing.

\begin{omfigure}
\begin{tikzpicture}[econ, x=1cm, y=0.24cm, font=\scriptsize]
  \node[anchor=east] (r) at (0,0) {data};
  \foreach \i/\lab in {0/outcome, 1/controls, 2/sample, 3/outliers}
    \node[font=\scriptsize\itshape] at (1.6*\i+1.3, 9.6) {\lab};
  \foreach \a in {-1,1} {
    \draw (0,0) -- (1.6,4*\a);
    \foreach \b in {-1,1} {
      \draw (1.6,4*\a) -- (3.2,4*\a+2*\b);
      \foreach \c in {-1,1} {
        \draw (3.2,4*\a+2*\b) -- (4.8,4*\a+2*\b+\c);
        \foreach \d in {-1,1} {
          \draw (4.8,4*\a+2*\b+\c) -- (6.4,4*\a+2*\b+\c+0.5*\d);
          \fill[omDef] (6.4,4*\a+2*\b+\c+0.5*\d) circle (1.2pt);
  }}}}
  \node[anchor=west, align=left] at (6.7,0) {16 paths:\\$2 \times 2 \times 2 \times 2$};
\end{tikzpicture}
\par\smallskip{\small Four binary choices --- which of two outcomes, whether to
control for age and schooling, whether to drop people over 50, whether to trim
outliers --- make sixteen analyses of the same data.}
\end{omfigure}

\omcode{code/bachelor-3/06-credible-empirical-work/listings/forking_paths.py}{1}{40}{Sixteen defensible analyses of data in which the treatment does nothing.}\label{lst:b3:credible-empirical-work:forking}

\omoutput{code/bachelor-3/06-credible-empirical-work/outputs/forking_paths.txt}{One dataset, then a thousand.}\label{out:b3:credible-empirical-work:forking}

In the first dataset no path is significant. Over a thousand datasets, the
first path --- the one a researcher could have written down in advance --- is
significant 6.4~\% of the time, close to the 5~\% it promises; at least one of the
sixteen paths is significant 21.4~\% of the time. % output: forking_paths
A researcher who reports the best path turns a 5~\% test into a 21~\% one.

\section{Low power and exaggerated estimates}\label{sec:b3:credible-empirical-work:power}

\begin{definition}[Exaggeration ratio; sign error]\label{def:b3:credible-empirical-work:exaggeration}
\index{exaggeration ratio}\index{sign error}
For a true effect $\beta \ne 0$ and a design with a given standard error, the
\emph{exaggeration ratio} (or type M error) is the expected absolute value of a
statistically significant estimate divided by $|\beta|$. A \emph{sign error} (or
type S error) is a statistically significant estimate whose sign is opposite to
that of $\beta$ (Gelman and Carlin, 2014).
\end{definition}

By \cref{prop:b3:credible-empirical-work:filter}, the \omterm{def:b3:credible-empirical-work:exaggeration}{exaggeration ratio} is
$1 + (s/\beta)\lambda(c - \beta/s)$, ignoring the small chance of a significant
estimate of the wrong sign. It depends only on power. The chart draws it.

\begin{omchart}{Exaggeration ratio of significant estimates against the power
of the study (normal approximation, test at 5~\%, two-sided, one tail ignored).}{Our
computation from the formula of
\cref{prop:b3:credible-empirical-work:filter}}
\begin{tikzpicture}
\begin{axis}[omtime, width=0.7\linewidth, height=0.38\linewidth, xmin=0, xmax=1,
  ymin=0, ymax=6, xlabel={Power}, ylabel={Exaggeration ratio},
  x tick label style={/pgf/number format/.cd, fixed, precision=1}]
  \addplot[omDef, thick, no marks, sharp plot] table[x=power, y=ratio] {figdata/bachelor-3/06-credible-empirical-work/out/exaggeration.csv};
  \draw[guide] (axis cs:0,1) -- (axis cs:1,1);
  \draw[guide] (axis cs:0.18,0) -- (axis cs:0.18,6);
  \node[font=\scriptsize, anchor=south west] at (axis cs:0.19,4.6) {power 0.18: ratio 2.4};
\end{axis}
\end{tikzpicture}
\end{omchart}

At a power of 0.8, the conventional target, significant estimates are 12~\% too
large on average; at 0.5, 41~\%; at 0.2, 125~\%. % figdata: exaggeration
\Cref{lst:b3:credible-empirical-work:filter} checks the formula by simulation, for
a treatment whose true effect is a tenth of a standard deviation and trials of
50 to 1{,}600 people per arm.

\omcode{code/bachelor-3/06-credible-empirical-work/listings/significance_filter.py}{1}{18}{The significance filter by simulation, against the formula (\texttt{filtered\_mean} in the chapter library).}\label{lst:b3:credible-empirical-work:filter}

\omoutput{code/bachelor-3/06-credible-empirical-work/outputs/significance_filter.txt}{Small trials publish large estimates.}\label{out:b3:credible-empirical-work:filter}

A trial with 50 people per arm has power 0.072 and, when significant, reports on
average an effect 4.85 times the truth; with 1{,}600 per arm, power 0.81 and a
ratio of 1.12. % output: significance_filter
Simulation and formula agree to within 0.04. At 50 per arm, \omterm{def:b3:credible-empirical-work:exaggeration}{sign errors} appear
too: a significant negative estimate has probability 0.007, against 0.072 for a
significant positive one, so about one significant result in eleven has the wrong
sign.

\begin{measure}[Power in economics]
Ioannidis, Stanley and Doucouliagos collected 64{,}076 estimates from more than
6{,}700 studies in 159 empirical literatures of economics. They estimate each
literature's true effect by the average of its adequately powered studies, and
each study's power against it. The median power is 18~\% or less, and in half of
the literatures nearly 90~\% of results are under-powered. Nearly 80~\% of the
reported effects are exaggerated, typically by a factor of two, and one third by
a factor of four or more. % ledger: b3ce:isd.power
At a power of 0.18 the formula above gives an \omterm{def:b3:credible-empirical-work:exaggeration}{exaggeration ratio} of 2.4 for a
significant result: the two methods agree in size.
\end{measure}

\section{The replication record of economics}\label{sec:b3:credible-empirical-work:replication}

\begin{definition}[Reproduction; replication]\label{def:b3:credible-empirical-work:replication}
\index{reproduction}\index{replication}
A \emph{reproduction} obtains a study's published results again from the same
data, using the authors' code or an independent rewriting of it. A
\emph{replication} repeats the study's design on new data --- a new sample from
the same population, or a new experiment --- and asks whether the result holds.
\end{definition}

A failed \omterm{def:b3:credible-empirical-work:replication}{reproduction} means an error or a missing step; a failed \omterm{def:b3:credible-empirical-work:replication}{replication}
may mean that the first result was a false positive, that it was exaggerated
by the filter, or that the effect differs between settings.

\begin{record}[Reproducing published papers from their code]
Chang and Li tried to reproduce 67 papers published in 13 economics journals,
using the data and code files the authors had provided. Without contacting the
authors they reproduced the key qualitative result of 22 papers (33~\%).
Leaving out 6 papers with confidential data and 2 that needed software they did
not have, they reproduced 29 of 59 (49~\%) with the authors' help.
% ledger: b3ce:cl.replicate
Most failures were missing files or code that did not run, not results that
were wrong; but a result that cannot be rerun cannot be checked.
\end{record}

\begin{record}[Replication of economics experiments]
Camerer and coauthors repeated 18 laboratory experiments published in two
leading economics journals between 2011 and 2014, each with enough participants
to have a power of at least 90~\% against the original effect. Eleven (61~\%)
found a significant effect in the original direction. On average the \omterm{def:b3:credible-empirical-work:replication}{replication}
effect was 65.9~\% of the original. % ledger: b2:camerer2016
Both numbers are what a literature with moderate exaggeration produces: most
effects real, most smaller than first reported.
\end{record}

\section{Pre-registration and pre-analysis plans}\label{sec:b3:credible-empirical-work:preregistration}

\begin{definition}[Pre-registration; pre-analysis plan]\label{def:b3:credible-empirical-work:pre-registration}
\index{pre-registration}\index{pre-analysis plan}
\emph{Pre-registration} records a study in a public, dated registry before its
data are collected or seen, with its design and main hypotheses. A
\emph{pre-analysis plan} is a document, registered in the same way, that fixes in
advance the outcomes, the sample, the specifications and the tests of the
study's main analysis.
\end{definition}

Registration attacks the filter: a registered study that is never reported is
visible as missing. A \omterm{def:b3:credible-empirical-work:pre-registration}{pre-analysis plan} attacks the forking paths: the first
path of \cref{out:b3:credible-empirical-work:forking}, chosen in advance, keeps
its 5~\% error rate. The American Economic Association decided in April 2012
to set up a registry for randomised trials in economics and the other social
sciences. % ledger: b3ce:aea.registry
The limits are real. Plans suit experiments, whose data do not exist before
the study; much of economics uses data that already exist, where no one can
check that the plan came first. And a plan that fixes everything leaves no room
for what the data teach; the usual answer is to report the planned analysis
first and label everything else exploratory.

\section{Meta-analysis and funnel plots}\label{sec:b3:credible-empirical-work:meta}

\begin{definition}[Meta-analysis; funnel plot]\label{def:b3:credible-empirical-work:meta-analysis}
\index{meta-analysis}\index{funnel plot}
A \emph{meta-analysis} combines the estimates of many studies of the same effect
into one. Its \emph{fixed-effect} mean weights each estimate by its precision,
$\bar\beta = \sum_i w_i\hat\beta_i / \sum_i w_i$ with $w_i = 1/s_i^2$, and has
standard error $1/\sqrt{\sum_i w_i}$; a \emph{random-effects} mean adds the
variance $\tau^2$ of true effects across studies to each $s_i^2$. A \emph{funnel
plot} draws each estimate against its precision $1/s_i$.
\end{definition}

Without selection, the funnel is symmetric: precise studies cluster near the
true effect at the top, imprecise ones scatter widely and evenly at the bottom.
With the filter, the bottom-left is missing --- small studies with small or
negative estimates were not published --- and the funnel leans.

\begin{omchart}{Funnel plot of a simulated literature: 400 studies of a true
effect of 0.1 (all dots), of which 138 are published (dark dots) when significant
positive results are always published and the others one time in ten.}{Simulated
by the authors (seed 400), the draws of
\cref{lst:b3:credible-empirical-work:funnel}}
\begin{tikzpicture}
\begin{axis}[omscatter, width=0.75\linewidth, height=0.42\linewidth, xmin=-0.6, xmax=0.8,
  ymin=0, ymax=34, xlabel={Estimate}, ylabel={Precision $1/s$}, legend pos=north east,
  x tick label style={/pgf/number format/.cd, fixed, precision=1}]
  \addplot[only marks, mark=o, mark size=1.4pt, draw=omDef!45,
    restrict expr to domain={\thisrow{published}}{0:0}]
    table[x=est, y=precision] {figdata/bachelor-3/06-credible-empirical-work/out/funnel.csv};
  \addplot[only marks, mark=*, mark size=1.4pt, fill=omDef, draw=omDef,
    restrict expr to domain={\thisrow{published}}{1:1}]
    table[x=est, y=precision] {figdata/bachelor-3/06-credible-empirical-work/out/funnel.csv};
  \draw[guide] (axis cs:0.1,0) -- (axis cs:0.1,34);
  \legend{not published, published}
\end{axis}
\end{tikzpicture}
\end{omchart}

\begin{definition}[Precision-effect test (FAT--PET)]\label{def:b3:credible-empirical-work:fat-pet}
\index{precision-effect test}\index{FAT--PET}
The \emph{FAT--PET} regression is the weighted least-squares regression of the
estimates on their standard errors, with weights $1/s_i^2$:
$\hat\beta_i = \beta_0 + \beta_1 s_i + e_i$. The funnel-asymmetry test (FAT)
tests $\beta_1 = 0$; the \emph{precision-effect test} (PET) tests $\beta_0 = 0$,
and $\hat\beta_0$, the predicted estimate of a study with no sampling error, is
the estimate of the effect corrected for selection.
\end{definition}

\begin{proposition}[Without selection, FAT--PET finds no asymmetry]\label{prop:b3:credible-empirical-work:fatpet}
Dividing the \omterm{def:b3:credible-empirical-work:fat-pet}{FAT--PET} regression by $s_i$ gives the ordinary regression $t_i =
\beta_1 + \beta_0 (1/s_i) + v_i$ of each study's $t$-statistic on its precision,
with $v_i = e_i/s_i$. If the estimates are unselected, $\hat\beta_i = \beta + s_i
u_i$ with $u_i \sim N(0, 1)$ independent of $s_i$, then this regression has
intercept $\beta_1 = 0$ and slope $\beta_0 = \beta$.
\end{proposition}

\begin{proof}
Weighted least squares with weights $1/s_i^2$ minimises $\sum_i (\hat\beta_i -
\beta_0 - \beta_1 s_i)^2/s_i^2 = \sum_i (t_i - \beta_0/s_i - \beta_1)^2$, which is
ordinary least squares of $t_i$ on a constant and $1/s_i$. Without selection,
$t_i = \beta (1/s_i) + u_i$, a linear model in $1/s_i$ with intercept 0, slope
$\beta$ and an error $u_i$ of mean zero independent of the regressor; OLS
estimates it without bias.
\end{proof}

Under the filter, imprecise studies are published only with large estimates,
so estimates rise with $s_i$ and $\hat\beta_1 > 0$. PET extrapolates the line to
$s = 0$, where no filter can act.

\omcode{code/bachelor-3/06-credible-empirical-work/listings/funnel.py}{1}{19}{\omterm{def:b3:credible-empirical-work:meta-analysis}{Meta-analysis} of the simulated literature, before and after the filter.}\label{lst:b3:credible-empirical-work:funnel}

\omoutput{code/bachelor-3/06-credible-empirical-work/outputs/funnel.txt}{The filter, seen and corrected.}\label{out:b3:credible-empirical-work:funnel}

With all 400 studies, every method finds about 0.1 and the FAT slope is
insignificant. Among the 138 published, the simple mean is 0.145, the
precision-weighted mean 0.127, and the FAT slope 0.762 is nearly four standard
errors from zero; PET gives 0.091, close to the truth. % output: funnel
PET is a model, not a cure: it assumes the bias is linear in the standard
error, and with few small studies, or when real effects differ by study size
(small studies of the most promising sites), it can mislead. The minimum-wage
literature of \cref{ch:b3:minimum-wage} is a real case.

\section{Reading a results table}\label{sec:b3:credible-empirical-work:table}

Most empirical results reach readers as a table. \Cref{lst:b3:credible-empirical-work:table}
prints one, from the job-training experiment of
\cref{ch:b3:difference-in-differences}.

\omcode{code/bachelor-3/06-credible-empirical-work/listings/reading_table.py}{1}{18}{Three nested regressions, printed as a journal table.}\label{lst:b3:credible-empirical-work:table}

\omoutput{code/bachelor-3/06-credible-empirical-work/outputs/reading_table.txt}{A results table: earnings in 1978, thousands of dollars.}\label{out:b3:credible-empirical-work:table}

Read it in this order. The columns are models; the rows are coefficients, each
with its standard error in parentheses below it (here robust ones, which the
table does not say --- a good table does). The stars mark $p$-values below 0.10
(*), 0.05 (**) and 0.01 (***). $N$ is the number of observations. Controls not
shown are listed in the code; a journal table lists them in a note. Then look at
what the stars hide. The training offer's coefficient falls from 1.794 to 1.651
and its standard error from 0.671 to 0.647: the estimate is all but unchanged,
yet column (3) has two stars where the others have three, because its
$t$-statistic, $1.651/0.647 = 2.55$, falls just short of 2.58. % output: reading_table
The stars draw lines through a continuous measure of evidence. Read the
coefficient and its standard error; the interval
$1.651 \pm 1.96 \times 0.647$, from 0.38 to 2.92 thousand dollars, says more than
any number of stars.

\section{The evidence scale of this volume}\label{sec:b3:credible-empirical-work:scale}

\begin{definition}[Evidence scale]\label{def:b3:credible-empirical-work:evidence-scale}
\index{evidence scale}
The \emph{evidence scale} of this volume grades the evidence for a causal claim:
\begin{enumerate}
\item several randomised or natural experiments that agree;
\item one strong design (a randomised trial, a convincing natural experiment), or
several that disagree in size;
\item correlations with controls and a plausible mechanism;
\item a theory with no design-based test yet.
\end{enumerate}
\end{definition}

\begin{measure}[Using the evidence scale]
Every \texttt{debate} box in this volume tags each study it cites with a grade.
A grade belongs to a claim, not to a paper: a randomised trial gives grade 2 for
the effect in its setting, and only agreement across settings gives grade 1.
\omterm{def:b3:credible-empirical-work:pre-registration}{Pre-registration}, a large sample and a reproduced analysis make a grade
trustworthy; a heap of significant small studies, an asymmetric funnel or a
failed \omterm{def:b3:credible-empirical-work:replication}{replication} lowers it. No grade turns a positive claim into a policy;
that needs values the evidence does not supply.
\end{measure}

\section{Exercises}\label{sec:b3:credible-empirical-work:exercises}

\begin{exercise}[$\star$]\label{exo:b3:credible-empirical-work:1}
In \cref{out:b3:credible-empirical-work:table}, which number is the standard
error of the training coefficient in column (2)? Give its $t$-statistic and the
approximate 95~\% interval. What would the table need to say for you to know
how the standard errors were computed?
\end{exercise}

\begin{exercise}[$\star$]\label{exo:b3:credible-empirical-work:2}
Compute $1 - 0.95^k$ for $k = 5$ and $k = 20$. Why is the figure for sixteen
correlated paths in \cref{out:b3:credible-empirical-work:forking} lower than the
formula's?
\end{exercise}

\begin{exercise}[$\star$]\label{exo:b3:credible-empirical-work:3}
A study's design gives a standard error of 0.02 for an effect whose true value
is 0.05. What is its power at 5~\% (ignore the wrong tail)?
\end{exercise}

\begin{exercise}[$\star$]\label{exo:b3:credible-empirical-work:4}
What would a \omterm{def:b3:credible-empirical-work:meta-analysis}{funnel plot} of an unbiased literature look like? Give one reason,
other than \omterm{def:b3:credible-empirical-work:publication-bias}{publication bias}, why a funnel might be asymmetric.
\end{exercise}

\begin{exercise}[$\star\star$]\label{exo:b3:credible-empirical-work:5}
Compute the \omterm{def:b3:credible-empirical-work:exaggeration}{exaggeration ratio} of a significant estimate at a power of 0.2,
using the normal approximation: find $\beta/s$, then $a$ and $\lambda(a)$, with
$\Phi^{-1}(0.8) = 0.8416$ and $\phi(0.8416) = 0.2800$.
\end{exercise}

\begin{exercise}[$\star\star$]\label{exo:b3:credible-empirical-work:6}
Twelve studies report estimates of the same effect, with standard errors (in
parentheses) of 0.02, 0.03, 0.04, 0.05, 0.06, 0.08, 0.10, 0.12, 0.15, 0.18, 0.20 and
0.25, and estimates 0.05, 0.07, 0.06, 0.11, 0.13, 0.17, 0.21, 0.26, 0.31, 0.38, 0.42 and
0.53 in the same order.
\begin{enumerate}
\item Compute the simple and the inverse-variance means.
\item Run \omterm{def:b3:credible-empirical-work:fat-pet}{FAT--PET}. What do you conclude about the effect?
\end{enumerate}
\end{exercise}

\begin{exercise}[$\star\star$]\label{exo:b3:credible-empirical-work:7}
How would a \omterm{def:b3:credible-empirical-work:pre-registration}{pre-analysis plan} change the result of
\cref{lst:b3:credible-empirical-work:forking}? What would it not change?
\end{exercise}

\begin{exercise}[$\star\star$]\label{exo:b3:credible-empirical-work:8}
Three studies estimate the same effect: 0.20 (s.e.\ 0.10), 0.10 (0.05) and 0.40
(0.20). Compute the fixed-effect mean and its standard error. Which study
dominates, and why?
\end{exercise}

\begin{exercise}[$\star\star$]\label{exo:b3:credible-empirical-work:9}
A well-powered \omterm{def:b3:credible-empirical-work:replication}{replication} of a published experiment finds an effect a third of
the original's size, not significant. List what this shows and what it does not
show about the original study.
\end{exercise}

\begin{exercise}[$\star\star\star$]\label{exo:b3:credible-empirical-work:10}
Under the conditions of \cref{prop:b3:credible-empirical-work:filter}, the
\omterm{def:b3:credible-empirical-work:exaggeration}{exaggeration ratio} is $1 + (s/\beta)\lambda(c - \beta/s)$.
\begin{enumerate}
\item Compute it at $\beta/s = 1$ and at $\beta/s = 3$.
\item Show that it tends to 1 as $\beta/s$ grows.
\item Now let estimates be published when $|\hat\beta/s| > 1.96$, of either sign.
At $\beta/s = 1$, compute the probability of publication and the share of
published estimates that are \omterm{def:b3:credible-empirical-work:exaggeration}{sign errors}.
\end{enumerate}
\end{exercise}

\begin{exercise}[$\star\star\star$]\label{exo:b3:credible-empirical-work:11}
Eight studies report 0.12 (0.05), 0.05 (0.04), 0.30 (0.10), 0.18 (0.06), $-0.02$
(0.05), 0.22 (0.08), 0.09 (0.03) and 0.15 (0.07).
\begin{enumerate}
\item Compute the fixed-effect mean and Cochran's $Q = \sum_i w_i (\hat\beta_i -
\bar\beta)^2$.
\item Estimate the between-study variance by DerSimonian and Laird, $\hat\tau^2 =
\max\{0, (Q - (k - 1))/(\sum w_i - \sum w_i^2/\sum w_i)\}$, and the
random-effects mean with weights $1/(s_i^2 + \hat\tau^2)$. Why is its standard
error larger?
\end{enumerate}
\end{exercise}

\begin{exercise}[$\star\star\star$]\label{exo:b3:credible-empirical-work:12}
Choose three published claims from earlier chapters of this volume --- one
from a randomised experiment, one from a natural experiment, one from
correlations --- and grade each on the \omterm{def:b3:credible-empirical-work:evidence-scale}{evidence scale}. Defend each grade, and
say what new study would raise it.
\end{exercise}

\section{Problem: Twenty Studies of a Training Programme}\label{sec:b3:credible-empirical-work:problem}

\begin{problem}[{Weekend problem --- twenty studies of a training programme}]\label{pb:b3:credible-empirical-work:twenty}
Brenland (invented) runs a job-training programme, and twenty randomised
evaluations of it have been published. Each reports the effect on trainees'
earnings a year later, in percent, with its standard error, and the number of
trainees in each arm. % figdata: studies
A summary of the literature, and a large pre-registered trial with 20{,}000
trainees per arm published later, are in the second table.
% figdata: studies_summary

\begin{center}\small
\pgfplotstabletypeset[col sep=comma, skip rows between index={10}{20},
  columns/study/.style={column name={Study}, string type},
  columns/n/.style={column name={$n$ per arm}, int detect},
  columns/est/.style={column name={Effect (\%)}, fixed, precision=1, zerofill},
  columns/se/.style={column name={s.e.}, fixed, precision=2, zerofill},
  every head row/.style={after row=\hline}, every last row/.style={after row=\hline}]
  {figdata/bachelor-3/06-credible-empirical-work/out/studies.csv}
\qquad
\pgfplotstabletypeset[col sep=comma, skip rows between index={0}{10},
  columns/study/.style={column name={Study}, string type},
  columns/n/.style={column name={$n$ per arm}, int detect},
  columns/est/.style={column name={Effect (\%)}, fixed, precision=1, zerofill},
  columns/se/.style={column name={s.e.}, fixed, precision=2, zerofill},
  every head row/.style={after row=\hline}, every last row/.style={after row=\hline}]
  {figdata/bachelor-3/06-credible-empirical-work/out/studies.csv}

\medskip
\pgfplotstabletypeset[col sep=comma,
  columns/stat/.style={column name={Statistic}, string type, column type=l},
  columns/est/.style={column name={Estimate}, fixed, precision=2, zerofill},
  columns/se/.style={column name={s.e.}, fixed, precision=2, zerofill, empty cells with={---}},
  every head row/.style={after row=\hline}, every last row/.style={after row=\hline}]
  {figdata/bachelor-3/06-credible-empirical-work/out/studies_summary.csv}
\end{center}

\medskip\noindent\textbf{Part I --- the naive average.}
\begin{enumerate}
\item How many of the twenty studies are significant at 5~\% ($z > 1.96$)? Check
the share in the summary table.
\item The simple mean is 6.52~\%. What does averaging the estimates with equal
weights assume about the studies?
\item Compute the inverse-variance weights of studies S16 and S04 and their
ratio.
\item Why is the inverse-variance mean, 3.78~\%, lower than the simple mean?
\item Give the 95~\% interval of the inverse-variance mean.
\item If the true effect were 3.78~\%, what power would a study the size of S04
have? What does that suggest about how S04 came to be published?
\end{enumerate}

\medskip\noindent\textbf{Part II --- funnel and FAT--PET.}
\begin{enumerate}[resume]
\item Sketch the \omterm{def:b3:credible-empirical-work:meta-analysis}{funnel plot} of the twenty studies, and say what it would look
like without selection.
\item Among the eight studies with standard errors above 5, how many estimates
lie above 3.78 and how many below? What does that show?
\item Interpret the FAT slope, 1.26 (s.e.\ 0.48), and its $t$-statistic.
\item Interpret the PET estimate, 0.93 (s.e.\ 1.23).
\item Why is PET read as the effect of a study with no sampling error?
\item If the true effect is 2~\%, compute the power of a study with a standard
error of 8 and the average size of its significant estimates. Compare with S02
and S04.
\end{enumerate}

\medskip\noindent\textbf{Part III --- a pre-registered trial.}
\begin{enumerate}[resume]
\item Give the trial's 95~\% interval.
\item Is the trial consistent with the inverse-variance mean? Compute the $z$
of their difference, treating the two as independent.
\item Is it consistent with PET?
\item What was the trial's power against a true effect of 2~\%?
\item Why does its registration matter for how its result is read?
\item Does the trial's estimate prove that PET was right?
\end{enumerate}

\medskip\noindent\textbf{Part IV --- the corrected estimate and its grade.}
\begin{enumerate}[resume]
\item Combine the PET estimate and the trial by inverse-variance weighting.
Give the combined estimate and its standard error.
\item Why should this combination be read with caution?
\item Grade the evidence for a positive effect on the scale of
\cref{def:b3:credible-empirical-work:evidence-scale}, before and after the trial.
\item A trainee would earn 20{,}000 dalens a year without training. Compute
the yearly gain at 6.52~\% and at 2.0~\%.
\item What does a reader of only the simple mean get wrong, and by how much?
\item State the results: the bias-corrected effect of the programme and its
grade.
\end{enumerate}
\end{problem}
